\ficha{Criação, 1906: os defensores}

\bloco{Fontes lidas}

\textit{Documentos Parlamentares}, volume de 1906. Li na íntegra os dois
discursos de David Campista, relator e autor do projeto, de 1º de setembro
(p.~233-269) e de 5 de outubro (p.~510-541). Dos demais defensores li a
abertura, o fecho e as passagens sobre taxa, padrão e fundos: Altino Arantes
(22 de agosto), Adolpho Gordo (3 de setembro), Rodolpho Paixão (19 de
setembro), Alberto Sarmento (20 de setembro), Barros Franco Júnior (5 de
outubro) e, no Senado, o parecer de Urbano Santos e a intervenção de Moniz
Freire (novembro).

\bloco{Síntese}

O relator não defende câmbio baixo. Defende a estabilidade como bem em si, a
uma taxa que descreve como a corrente, e apresenta a Caixa como complemento, e
não como negação, da política de Joaquim Murtinho. A divisão de funções que ele
propõe é simples: a Caixa contém a alta, os fundos de 1899 contêm a baixa. A
bancada paulista acrescenta o argumento do café e aceita 15 como ponto de
partida. No Senado, o parecer nega que a estabilidade a 15 seja definitiva.

\bloco{David Campista}

A estabilidade vale mais que o nível:

\cit[P:1:242]{Um governo que, durante quatro annos, conseguisse manter a mesma
taxa cambial, embora relativamente baixa, seria, a meu ver, mais benemerito do
que aquelle que, encontrando um cambio de 10, o elevasse no mesmo periodo a
25}{\dpa{p.~236}}

A lavoura, diz ele, \tr[P:1:274]{pedia tão sómente a estabilidade como condição
de vida} \dpa{p.~268}. A taxa de 15 é escolhida porque \tr[P:1:541]{tem sido,
mais ou menos, a taxa corrente} \dpa{p.~535}.

A divisão de funções entre a Caixa e os fundos de 1899:

\cit[P:1:264]{a lei em projecto é feita para conter a alta e os fundos de
resgate e garantia para conter a baixa}{\dpa{p.~258}}

Sobre o padrão, Campista só admite a quebra \tr[P:1:247]{desde que seja seguida
da immediata conversão do meio circulante} \dpa{p.~241}, o que o projeto não
faz. Sobre a herança de 1899, afirma que \tr[P:1:538]{o projecto é o
complemento dessa politica} \dpa{p.~532}.

Ele também localiza a resistência ao projeto:

\cit[P:1:521]{As resistencias ficaram na Capital, subindo das agitações
interessadas da rua da Alfandega ás regiões do officialismo, com escalas pela
imprensa local}{\dpa{p.~515}}

Na mesma passagem ele trata da imprensa do Rio. Diz que seguiu os conselhos
do \textit{Correio da Manhã}, cujo artigo de 13 de janeiro de 1906 a favor da
fixação do câmbio transcreve, e aponta a mudança de \textit{O Paiz}:

\cit[P:1:521]{Foi, por exemplo, no O Paiz, que, ainda este anno, encontrei as
melhores defesas das idéas pelas quaes me bato agora}{\dpa{p.~515}}

\cit[P:1:521]{E é agora o O Paiz que, como recentemente se viu, chega ao
extremo de appellar para os meios violentos como protesto contra a
reforma}{\dpa{p.~515}}

É testemunho de parte interessada, mas datado de 5 de outubro de 1906 e
coerente com o que a ficha 8 registra nas páginas do próprio jornal.

\bloco{A bancada paulista e os demais}

Adolpho Gordo define o objetivo como \tr[P:1:296]{fixar o cambio em uma taxa
que reflicta a nossa situação economica} \dpa{p.~290} e adverte que a alta
brusca \tr[P:1:309]{Será a ruina da nossa lavoura} \dpa{p.~303}. Altino Arantes
apresenta a Caixa como passo para \tr[P:1:137]{o advento da circulação
metallica no Brazil} \dpa{p.~131}. Rodolpho Paixão considera a
\tr[P:1:375]{taxa de 15 pence por mil réis, que eu considero muito acceitavel}
\dpa{p.~369}. Alberto Sarmento emenda o projeto para que a taxa só possa ser
elevada \tr[P:1:393]{por lei do Congresso} \dpa{p.~387}. Barros Franco Júnior
liga a reforma ao Convênio de Taubaté: o país \tr[P:1:515]{em 1906, vae
conseguir a sua emancipação economica} \dpa{p.~509}.

No Senado, o parecer de Urbano Santos sustenta que o projeto preserva o padrão
legal: \tr[P:1:565]{Não é, pois, uma medida definitiva a estabilidade}
\dpa{p.~559}. Moniz Freire vota a favor e diz que suas emendas
\tr[P:1:602]{vêm em apoio do plano} \dpa{p.~596}. Rui Barbosa vota
\textit{sim} na chamada de 26 de novembro (p.~602) e não discursa sobre a
Caixa no volume.

\bloco{Achados}

\begin{enumerate}[leftmargin=*,nosep]
\item Os defensores da Caixa não se descrevem como partidários do câmbio baixo.
Codificar a defesa da Caixa como defesa do câmbio baixo seria erro de
construto: o que separa Campista dos valorizadores é o ritmo da alta e a
garantia metálica, não o destino declarado.
\item Campista e seus críticos disputam a herança de Murtinho. O relator se diz
continuador; Serzedelo e Barbosa Lima lhe negam esse título (ficha 1).
\item A emenda Sarmento, que reserva ao Congresso a elevação da taxa, é a
origem do debate de 1910 sobre competência (ficha 5).
\end{enumerate}

\bloco{Limites}

Só os dois discursos de Campista foram lidos por inteiro. A posição de Rui
Barbosa em 1906 está documentada pelo voto; a ficha 3 registra o que a imprensa
lhe atribuiu antes dele.
